\section{Introduction}
\label{sec:intro}

This paper studies the bin-packing problem.
Here we are given a set of $n$ items. The $i\Th$ item has size $s_i \in (0, 1]$.
Our goal is to pack the items into a small number of bins.
Formally, we must partition the items into sets $B_1, \ldots, B_m$ such that
$m$ is small and for each bin $j \in [m]$, we have
\[ \sum_{i \in B_j} s_i \le 1. \]

The \emph{harmonic algorithm} \cite{lee1985simple} is a simple online algorithm for this problem.
This algorithm has an integer parameter $M \ge 1$.
For each item of size $x \in (0, 1]$, define its \emph{category} $\categ_M(x)$
as $\min(M, \floor{1/x})$, and its \emph{harmonic weight} $w_M(x)$ as
\[ w_M(x) \defeq \begin{cases}
    \frac{1}{\floor{1/x}} & \text{ if } 1/M < x \le 1
    \\ \frac{M}{M-1}x & \text{ if } 0 < x \le 1/M
\end{cases}. \]
The harmonic algorithm packs all items of the same category using the Next-Fit algorithm.

\subsection{Weight Functions}

The analysis of several bin-packing algorithms follow a technique called \emph{weight functions}.
The key idea is to devise a function $w: (0, 1] \to \mathbb{R}_{> 0}$
and find constants $\alpha, \beta \in \mathbb{R}_{\ge 0}$ having the following properties:
\begin{enumerate}
\item The algorithm packs all items into less than $\sum_{i=1}^n w(s_i) + \beta$ bins.
\item If a subset $S$ of items can be packed into a bin, their total weight is at most $\alpha$. Formally,
    \[ \sum_{i \in S} s_i \le 1 \implies \sum_{i \in S} w(s_i) \le \alpha. \]
\end{enumerate}
Then it can be shown that the algorithm packs all items into less than $\alpha m^* + \beta$ bins,
where $m^*$ is the optimal number of bins.
Here $\alpha$ is called the \emph{asymptotic approximation factor}.

The harmonic algorithm uses the harmonic weight function $w_M$,
and it can be shown that it packs all items into less than $\sum_{i=1}^n w_M(s_i) + (M-1)$ bins.
In this paper, we show an upper-bound on the total weight of items that can be packed into a bin.

\subsection{Statement of this Paper's Main Result}

We first define $Q_M$, the supposed asymptotic approximation factor of the harmonic algorithm.

For any $n \in \mathbb{N}$, let $k_n$ be the $n\Th$ \emph{modified Sylvester number} \cite{oeisA007018},
defined recursively as follows:
\[ k_n \defeq \begin{cases}
    1 & \text{ if } n = 1
    \\ k_{n-1}(k_{n-1}+1) & \text{ if } n > 1
\end{cases}. \]
For any $n \defeq \mathbb{N}$, define $T_n$ to be the reciprocal of the first $n$ modified Sylvester numbers:
\[ T_n \defeq \sum_{i=1}^n \frac{1}{k_i}. \]

For any number $n \in \mathbb{N}$, define its \emph{modified Sylvester inverse} to be the
largest number $i \in \mathbb{N}$ such that $n \le k_i$.
(If $n \ge 2$, then we get that $i \ge 2$ and $k_{i-1} < n \le k_i$.)

For any integer $M \ge 2$ and its modified Sylvester inverse $i$, define $Q_M$ as
\[ T_i + \frac{1}{M-1}\frac{1}{k_i}. \]
Define $T_{\infty}$ and $Q_{\infty}$ as
\begin{align*}
T_{\infty} &\defeq \lim_{n \to \infty} T_n
& Q_M &\defeq \lim_{M \to \infty} Q_M
\end{align*}
Then $T_{\infty} = Q_{\infty} \approx 1.691$.
In practice, for $M \ge 12$, $Q_M \approx Q_{\infty}$.

We now state the main result that we prove.

\begin{restatable}{theorem}{thmHarmonicWeightBound}
\label{thm:harmonic-bound}
For all $n, M \in \mathbb{N}$ and $y \in (0, 1]^n$ such that
$M \ge 2$ and $\sum_{i=1}^n y_i \le 1$, we get
\[ \sum_{i=1}^n w_M(y_i) \le Q_M. \]
\end{restatable}
